\newcommand{\CaseID}{walker\_2019}
\newcommand{\CaseTitle}{2019 Walker fire reconstruction against GOFER observations}
\newcommand{\FireName}{Walker}
\newcommand{\FireYear}{2019}
\newcommand{\EventStart}{4 September 2019 at 07:00~UTC}
\newcommand{\EventEnd}{12 September 2019 at 22:00~UTC}
\newcommand{\DurationHours}{207}
\newcommand{\CRSName}{UTM zone 10N (EPSG:32610)}
\newcommand{\FuelVersionText}{LANDFIRE release 2.0.0\_2019}
\newcommand{\FireSummary}{The supplied GOFER event record contains 193 hourly perimeters, a final area of 267.693~km$^2$ and an official size of 54,612~ac. The configuration uses 1 timed ignition point(s).}
\newcommand{\ConfigurationText}{The supplied simulation runs for 193~h, with an initial timestep of 1~s, target Courant--Friedrichs--Lewy number of 0.2, 194 hourly weather fields and 1 ignition point(s). Arrival time, spread rate, fireline intensity and fire-size statistics are requested at the supplied dump interval. Crown-fire model 1 and firebrand spotting are enabled. Unspecified settings retain the inspected source defaults, including fixed random seed 2024, foliar moisture 90\% and maximum timestep 600~s. A single simulation member is configured.}
\newcommand{\WeatherOffsetText}{Simulation and weather time zero is 14~h after the GOES ignition hour; comparisons apply this offset explicitly.}
\newcommand{\ElmfireRevision}{8981f14e0170}
\newcommand{\ElmfireBranch}{yiren}
